Martinez\documentclass[11pt]{article}

\usepackage{fullpage}
\usepackage{graphicx}
\usepackage{hyperref}
\usepackage{amsmath}
\usepackage{float}
\usepackage{xurl}

\title{A Tool for Signal Chain Simulation in Particle Physics}
\author{Octavio Gamez and Desiree Martinez}
\date{July 31st, 2026}

\begin{document}

\maketitle

\begin{abstract}

In particle physics, signal processing is very important in the electronics chain and for the pre-planning phases of an experiment. However, by going from detector to digitizer, experimenters can use a full SPICE simulator to get their signal, but doing so requires a detailed schematic of electrical components, which is impractical during the early design phase of an experiment. Our solution to this is an open source tool that experimenters can use to plan an analog signal to the digitizer path while going through specified components to see how the waveform distorts from the original signal that was produced.
  
\end{abstract}

\section{Introduction}

Planning the analog signal path isn't thought to be as important of a step when doing signal processing. Experimenters usually start from the detector (photomultiplier tube, drift chamber, or pixel detectors such as GEMS), and give a brief thought of how the signal will travel from the source to the digitizer. The signal would travel through cables, splitters, amplifiers, and discriminators and would be built based on previous experiments without truly understanding the effect it could have on the signal. Using this approach does work well for those familiar with the hardware and setups they know well, but fails to anticipate problems such as reflections that can corrupt timing, signal loss due to long cables, or distortion that sneakily messes up a fast pulse before these problems show up in real data.

There are tools to analyze the path of the analog signal to the digitizer with high precision. Experimenters can use a full SPICE simulator that can accurately find the distortions in the signal from the different components it runs through. But getting this level of accuracy requires a detailed schematic of all the electrical components that are involved, which would be unrealistic for the pre-planning phase of an experiment. Because of this, experimenters are left between choosing intuition from past experiments or building a complete circuit simulation to get their results.

This project is to fix that gap with an open source tool that is designed specifically for the pre-planning stage of the experiment. This tool is used to let an experimenter describe their signal generator and use it through different components to the digitizer. Depending on what the experimenter chooses, the signal to run through our tool will calculate the predicted outcome of the final result of the signal. This tool will help identify whether reflections occur along with how large they will be, how much the waveform will distort by the time it reaches the digitizer, and how much signal amplitude could be taken or added to the original signal. The goal of this tool is to build something simple enough to use during early planning but also sturdy enough for experimenters to know what they want before building anything.
%-------------------------------------------------------------
\section{Background}

To understand how our tool can model a signal path to the digitizer, it helps to understand the physical behavior of each piece through which the signal will be filtered. This can include the shape of the signal at its source, what happens as it moves through cables and components, to how it is eventually timed and measured. The following sections will talk more about these concepts.

% Check the math and equations to make sure it all makes sense
\subsection{PMT Signal}
In short, a photomultiplier tube (PMT) is a device used to amplify low levels of light. A signal produced by a PMT may last anywhere from 1 to 50 ns. Photons strike at one end of the tube, which holds a photocathode. The photocathode then produces photoelectrons via the photoelectric effect; these are then further amplified by the successive dynodes before being collected by the anode in the form of current caused by the electron avalanche. When the current transits through a load resistor, a signal of voltage as a function of time is extracted. The pulse is fairly predictable: It has a rapid rise before experiencing slow decay. This shape is approximated with a bi-exponential function. There are many forms of the bi-exponential, one being:

$$
f(t) = A \cdot \left[ e^{-\frac{t}{\tau_{\text{fall}}}} - e^{-\frac{t}{\tau_{\text{rise}}}} \right]
$$


Where:
\begin{align*}
f(t) &\quad\to \text{output values at a specific time } (t) \\
A    &\quad\to \text{an arbitrary scaling factor} \\
t    &\quad\to \text{time} \\
\tau_{\text{rise}} &\quad\to \text{Time constant controlling the signal rise time}\\
\tau_{\text{fall}} &\quad\to \text{Time constant controlling the signal fall time}
\end{align*}

The two time constants, ${\tau_{\text{rise}}}$ and ${\tau_{\text{fall}}}$, can be adjusted to match a specific PMT. Real pulses vary in amplitude from one event to the next even under identical conditions.

\subsection{Transmission Lines and Impedance matching}
Characteristic impedance is defined as the ratio of voltage to current for a wave traveling through a transmission line. It is given by:

$$Z_0 = \sqrt{\frac{R + j\omega L}{G + j\omega C}}$$


A mismatch between the characteristic impedance and the load and source impedance along the circuit pathway will cause a signal to reflect. Signal reflection can be quantified by the reflection coefficient:

$$
\Gamma = \frac{(Z_{\text{load}} - Z_0)}{(Z_{\text{load}} + Z_0)}
$$

Where:
\begin{align*}
\Gamma &\quad\to \text{reflection coefficient} \\
Z_{\text{load}}    &\quad\to \text{impedance of the load} \\
Z_0    &\quad\to \text{characteristic impedance} \\
\end{align*}

 A matched connection ($Z_0$ = $Z_{\text{load}}$) gives no reflection ($\Gamma$ = 0). An open or short end cable yields:

$$
\Gamma_{\text{open}} = \lim_{Z_{\text{load}} \to \infty} \frac{Z_{\text{load}} - Z_0}{Z_{\text{load}} + Z_0} = +1
$$

for an open cable, and: 

$$
\Gamma_{\text{short}} =  \frac{0- Z_0}{0 + Z_0} = -1
$$

for a short cable. A reflection coefficient of 1 indicates that the entire signal was reflected. A reflection coefficient of -1 indicates that the entire signal was reflected and experienced an inversion in its polarity .

Because a signal moves through a transmission line at a finite speed, propagation delay is introduced to the signal. Delay is given by:

$$
t_{delay} = \frac{L}{v_f \cdot c}
$$

Where:
\begin{align*}
t_{delay} &\quad\to \text{the signal's shift in seconds} \\
L   &\quad\to \text{length of the cable in meters} \\
v_f    &\quad\to \text{velocity factor (ratio of transmit speed to the speed of light)} \\
c    &\quad\to \text{the speed of light)} \\
\end{align*}

Using the time delay, a delayed signal can be expressed as:

$$
V_{delayed}(t) = V_{orginal}(t - t_{delay})
$$

\subsection{Signal Integrity and Decay}

Travel through a cable causes a signal to undergo attenuation. It is common to have cables that measureMartinez attenuation in decibels per meter. In order to convert this into a usable scalar we first multiply this quantity by the cable's length:

$$
A_{dB} = \alpha \cdot L
$$

Where:
\begin{align*}
A_{dB} &\quad\to \text{the attenuation experienced by the signal in decibels} \\
\alpha   &\quad\to \text{attenuation per meter in decibels} \\
L   &\quad\to \text{length of the cable in meters} \\
\end{align*}

We can convert this into a usable, unitless, coefficient through base ten exponentiation: 

$$

$$


\subsection{Timing Extraction}

Somewhere in most signals, an analog pulse needs to be turned into a single moment that marks when a signal “arrived”. The simplest method is the leading/trailing edge discriminator, it does this by marking the moment the signal crosses a fixed voltage. It works, but a flaw is that larger pulses cross that point earlier than smaller pulses even if they both start at the same time. This is called a timing walk. But a constant fraction discriminator fixes this by triggering a fixed percentage of each pulses own height instead of the fixed voltage, that way the timing stays consistent no matter the pulse size.

\subsection{Existing Modeling Tools}
% MIGHT get rid of this subsection completely im not sure if its neccesary
Circuit simulation tools already exist. SPICE-based simulators can model reflections, loss, and distortion. However, SPICE needs a fully detailed schematic like exact components and values, which is not ideal during the early phases of planning
%-------------------------------------------------------------
\section{Methods/Process}
In this project, we only have a set amount of components that the signal would run through before it reaches the digitizer: splitter, amplifier, leading/trailing edge discriminator, constant fraction discriminator, cables, connectors, and terminators.
So when doing this project, we both worked on making classes in python for each of the components as well as making a signal generator (PMT Signal), circuit analyzer, and digitizer.

\begin{figure}[!htbp]
    \centering
    \includegraphics[width=0.7\textwidth]{ComponentsClasses_Screenshot.png}
    \caption{Component Classes}
\end{figure}

\subsection{Amplifier}
The amplifier class takes in a signal, subtracts it's baseline, and scales it by a gain factor. Low/high frequency cutoffs can also be applied by limiting how fast the voltage can change.
\subsection{Cable}
The cable class does multiple things at once in its code, it calculates how much the signal weakens while traveling, calculates the reflections that happen at the beginning of the cable and the end, and how much of the signal actually enters the cable.
\subsection{Circuit Analyzer}
The circuit analyzer class is the component responsible for measuring and comparing every stage of the signal chain and gathering the result into a final report. It does not modify the signal but instead observes each waveform after it goes through a component, it will store the results and after multiple tests have been ran it can compare any two of them and produce a summary of how the signal changed across the chose components.
\subsection{Connector/Terminators}
Connectors are junction points, anywhere impedance can change and parts of the signal can reflect. Terminators represent what sits at the end of the line and is defined by a single impedance value. Both of these rely on the reflection coefficient.
\subsection{Discriminators}
Leading/trailing edge discriminator marks a timing point when the signal crosses a fixed voltage. Like said previously, this brings up the timing walk where the marked time shifts depending on the pulse amplitude. The constant fraction discriminator is intended to correct this by triggering on a fixed fraction of each pulses own peak instead.
\subsection{Digitizer}
%Can try to make it more spepcific, kinda hard when all it does is just digitizze what the signal looks like
Our digitizer code models an analog to digital converter that turns the PMT signal into a digital representation of what the signal would look like. Simply put, it does everything a real digitizer would show about a signal.
\subsection{Signal Generator}
The signal generator simulates a synthetic PMT waveform, a fast rise followed with a slow decay. The user can control the amplitude, rise and fall time. This was the first class we built because every other component that the signal would go through before reaching the digitizer needed something to act on.
\subsection{Splitter}
The splitter is modeled as a three-resistor network. The splitter only uses resistors and because of this it lowers the voltage without changing the signals shape or delaying the signal, and each output retrieves exactly half of the input voltage.
%-------------------------------------------------------------
\section{User-Guide/Results}
For the user-guide this is the link to our GitHub, it will have all the steps to follow and code to read when using this tool.
\href{https://github.com/AverageAtBest000/SignalProcessingSimulator.git}{Octavio Gamez and Desiree Martinez: Signal Processing GitHub}.



\Large
\textbf{Results}

\normalsize

\begin{figure}[!htbp]
    \centering
    \includegraphics[width=0.7\textwidth]{Results_SignalProcessing.png}
    \caption{Results}
\end{figure}

As in the figure shown above, this should be what you see when you first run our code. This code shows the output for our PMT signal passing through a passive splitter to the digitizer. The original signal is roughly 0.046 in amplitude when the splitter comes into the picture both channels fall to about half of that value, around -0.023, at their peak. The two splitter channels overlap, which is expected, confirms that the splitter produces two identical outputs. The digitized version of split channel 2 drops by roughly half, making it roughly 0.0112. Even going through the splitter on all four of the signal/channels, the pulse’s overall timing and shape are preserved, along with the peaks occurring at approximately the same point in each case. Overall, this graph shows the splitter function and digitizer do affect the amplitude but also do not distort the pulse’s timing.

%-------------------------------------------------------------
\section{Limitations}
% Add more limitations and next steps
% In case you dont know how to just click enter for the next line and use the backslash \ and type item (\item)

Because this tool was built in a short time frame (four weeks) and was focused around the core signal path rather than every possible real-world detail, it leaves room for limitations to fully unlocking the potential the tool can have.
Some limitations that could come from what we have for our tool now are:
%add more stuff im blanking on most of the stuff my head hurts
\begin{itemize}
    \item We only had four weeks time to build this
    \item Noise is a rough estimate
    \item Tool stops at digitizer 
    \item Reflection only bounces back once (doesn't account for the signal bouncing back and fourth multiple times)
    \item Has not been tested against real hardware
\end{itemize}
%-------------------------------------------------------------
\section{Next Steps}

Most of the limitations aren't dead ends that should be tossed away but are a list of ideas that haven't been finished yet or what could be improved in the future.
As we continue to work on this project and look towards the future, some ideas that we could implement are as follows:
\begin{itemize}
    \item Test the tool against real measurements to check accuracy and compare
    \item Get it to support multiple reflections bouncing back and fourth
    \item Add more components and source types
    \item enhancing different features that are already implemented in our design
\end{itemize}
%-------------------------------------------------------------
\section{Conclusion}

Planning the analog signal path for particle physics has always been between choosing two paths: experience-based estimations and building a fully detailed simulation. These paths are not suited for the early planning stages of an experiment. This project was built with the idea to add another path for experimenters to use for their pre-planning stages. It's set to be an open-source tool that has a source, components that could mess up a signal, and giving back results about what experimenters want to know (reflections, distortion, and signal loss) before any hardware is built.

%-------------------------------------------------------------
\section{Acknowledgment}

We would like to thank our mentor Juan Carlos Cornejo for his help in guiding us along the journey of creating this tool and giving his valuable feedback. We would like to thank Rhonda Crespo and Mark Galassi for letting us take on this internship and learning advanced computing techniques along with expanding our research skills. We would also like to acknowledge our co-workers/friends that we met through this internship, we all learned and grew together in the four weeks of this internship.
% Might get rid of the last sentence it doesnt seem all that relevent, TBD

%-------------------------------------------------------------
\newpage
\section{References}

\begin{itemize}
\small
    \item AverageAtBest. “GitHub - AverageAtBest000/SignalProcessingSimulator.” GitHub, \url{github.com/AverageAtBest000/SignalProcessingSimulator.git}.
    \item Cornejo, Juan Carlos. “Signal Processing Summer Project Ideas.” Signal Processing Summer Project Ideas, project idea, 10 July 2026, p. Abstract.
    \item GeeksforGeeks. “Spectrum Analysis in Python.” GeeksforGeeks, 23 July 2025, www.geeksforgeeks.org/artificial-intelligence/spectrum-analysis-in-python.
    \item Journal of Open Source Software. joss.theoj.org.
\end{itemize}

\end{document}